\section{Literal analytic standard monomials over the centre base}
\label{sec:hironaka-actual-standard-normal-bases}
Let $B=\mathcal O_{\mathbb C^d,0}$, $R=\mathcal O_{\mathbb C^{d+t},0}$,
$N=(z_1,\ldots,z_t)$, $M=\mathfrak m_B R$, and $T:R\to\mathbb C[[u,z]]$ be the actual Taylor map.
For an arbitrary original ideal $I\subset R$, write $K_f=T(I+M)\mathbb C[[u,z]]$.
The finite index $S_q$ consists of the exponents of total degree $<q$ outside the genuine degree-first diagram of $K_f$.
All monomial lifts below are the literal analytic coordinate products, rather than unspecified lifts of residue classes.

\begin{theorem}[Prescribed lifts of the actual complex residue basis]
\label{hir:actual-prescribed-complex-basis-lift}
\uses{hir:actual-base-residue-field,hir:actual-base-residue-tensor}
\lean{Hironaka.HironakaActualStandardNormalBases.exists_basis_of_complex_basis}
\leanok
If $Q$ is finitely presented and flat over $B$, and a prescribed family $v_i\in Q$ has
$1\otimes v_i$ as a complex basis of $\mathbb C\otimes_BQ$, then there is a $B$-basis $b$ with $b_i=v_i$ for every $i$.
\end{theorem}
\begin{proof}
Origin evaluation gives the proved field equivalence $e:\kappa(B)\simeq\mathbb C$.
The genuine tensor equivalence $\kappa(B)\otimes_BQ\to\mathbb C\otimes_BQ$ sends $a\otimes v$ to $e(a)\otimes v$;
it is linear after restricting the complex action through $e$.
Thus the displayed complex basis transports to the literal family $1\otimes v_i$ over $\kappa(B)$.
Flatness makes $\mathfrak m_B\otimes_BQ\to Q$ injective.
The local finite-presentation basis-lifting theorem applies to this prescribed family and gives $b_i=v_i$.
In particular it does not replace the displayed lifts by other elements with the same residues.
\end{proof}

\begin{theorem}[Actual complex fibre standard monomial basis]
\label{hir:actual-normal-fibre-standard-basis}
\uses{hir:relative-normal-original-fibre-tensor-algebra,hir:relative-normal-original-fibre-maximal-power,hir:actual-original-formal-whole-jet-algebra-equivalence,hir:actual-formal-standard-monomial-whole-jet-basis}
\lean{Hironaka.HironakaActualStandardNormalBases.normalFibreFormalJetEquiv,Hironaka.HironakaActualStandardNormalBases.normalFibreFormalJetEquiv_tmul_mk,Hironaka.HironakaActualStandardNormalBases.complexStandardMonomialBasis,Hironaka.HironakaActualStandardNormalBases.complexStandardMonomialBasis_apply}
\leanok
For every $I,q$, the actual algebra maps identify
$\mathbb C\otimes_B R/(I+N^q)$ with $\mathbb C[[u,z]]/(K_f+(u,z)^q)$, preserving each numerator as its actual Taylor series.
The literal analytic monomials indexed by $S_q$ therefore give a complex basis of this residue tensor.
\end{theorem}
\begin{proof}
The whole actual fibre tensor equivalence has target $R/(I+M+N^q)$.
The proved ideal identity $I+M+N^q=I+M+\mathfrak m_R^q$ identifies this with the whole point quotient of the same ideal $I+M$.
The genuine original/formal whole-jet equivalence then gives the stated formal quotient, sending $[g]$ to $[Tg]$.
The existing formal standard-monomial equivalence has basis vectors $[u^\beta z^\gamma]$.
Its pullback through these actual equivalences is $1\otimes[u^\beta z^\gamma]$, since Taylor sends each analytic monomial to that same formal monomial.
No freeness of the original $B$-module is used in this residue assertion, including $q=0$.
\end{proof}

\begin{theorem}[The standard labels are genuine normal labels]
\label{hir:actual-fibre-standard-normal-labels}
\uses{hir:actual-normal-fibre-standard-basis}
\lean{Hironaka.HironakaActualStandardNormalBases.tangentVariable_mem_fibreFormalIdeal,Hironaka.HironakaActualStandardNormalBases.standardIndex_tangent_zero,Hironaka.HironakaActualStandardNormalBases.analyticMonomial_standardIndex_eq_normal,Hironaka.HironakaActualStandardNormalBases.normalStandardIndex,Hironaka.HironakaActualStandardNormalBases.normalStandardIndexEquiv}
\leanok
Every $\alpha\in S_q$ has zero tangent block.
Consequently $S_q$ is canonically equivalent to the normal exponents $\gamma$ of degree $<q$ with $(0,\gamma)\notin\Gamma(K_f)$, and its analytic monomials are literally the original normal coordinate products $z^\gamma$.
\end{theorem}
\begin{proof}
Each tangent coordinate $u_i$ is the actual tangent pullback of an element of $\mathfrak m_B$, so its Taylor variable lies in $K_f$.
Its exponent $e_i$ is in the actual diagram. If $\alpha_i>0$, upperness puts $\alpha$ in that diagram, contradicting standardness.
Thus all tangent exponents vanish. Splitting the two actual exponent blocks identifies $\alpha$ with $(0,\gamma)$; the total degree then equals the normal degree.
These explicit splitting and joining maps are inverse and preserve the original monomial products, including empty coordinate blocks and $q=0$.
\end{proof}

\begin{theorem}[Literal analytic standard monomial bases]
\label{hir:actual-normal-truncation-standard-basis}
\uses{hir:actual-prescribed-complex-basis-lift,hir:actual-normal-fibre-standard-basis,hir:actual-fibre-standard-normal-labels}
\lean{Hironaka.HironakaActualStandardNormalBases.exists_standardMonomialBasis,Hironaka.HironakaActualStandardNormalBases.exists_normalStandardMonomialBasis}
\leanok
If the whole actual normal truncation $R/(I+N^q)$ is finite free over $B$, then it has a $B$-basis indexed by $S_q$, whose vectors are the literal analytic monomial classes. Reindexing by the proved normal labels gives precisely the classes $[z^\gamma]$ as a $B$-basis.
\end{theorem}
\begin{proof}
Finite freeness implies finite presentation and flatness.
Apply the prescribed-lift theorem to the literal analytic monomial family and its proved complex residue basis.
Every resulting basis vector equals exactly the specified analytic quotient class. Reindex with the preceding explicit normal-label equivalence; the monomial equality gives the precise normal coordinate product at every label.
\end{proof}

\begin{theorem}[Original constant-HS centre supplies all literal bases]
\label{hir:actual-original-centre-standard-monomial-bases}
\uses{hir:actual-original-centre-normal-layers-finite-free,hir:actual-normal-truncation-standard-basis}
\lean{Hironaka.HironakaActualStandardNormalBases.exists_original_hs_constant_centre_standard_monomial_bases}
\leanok
The original coherent ideal and regular constant-Hilbert--Samuel centre admit one genuine flattening $\Phi$ and the same origin ideal $I_0$ such that all $R/(I_0+N^q)$ have the preceding literal bases.
\end{theorem}
\begin{proof}
Choose the single flattening and original origin ideal from the proved all-degree original normal-layer freeness producer.
The actual short exact truncation sequences and actual zero quotient give all normal truncations free by induction; they are finite since they are quotients of the genuine finite normal-monomial ambient truncation.
For each $q$ apply the preceding theorem to that same $I_0$.
No second existential chart, flatness, basis, initial-diagram, or compatibility certificate is supplied.
\end{proof}

This produces the actual fibre-diagram monomial bases.
Comparison of this diagram with the original ideal diagram, arbitrary-ideal relative formal quotient transport, and the whole reduced-vertex normal-power assertion remain subsequent results; no Bennett or general G8 closure follows from the basis statement alone.
